\documentclass[11pt]{article}
\usepackage[margin=1in]{geometry}
\usepackage[T1]{fontenc}
\usepackage{graphicx}
\usepackage{hyperref}
\usepackage{enumitem}
\setlength{\parskip}{0.5em}
\setlength{\parindent}{0pt}

\title{state space models for sequence modeling}
\author{arXiv GraphRAG --- phase-6 deep research}
\date{\today}

\begin{document}
\maketitle

\begin{abstract}
An outstanding review that masterfully connects historical bottlenecks to the current frontier of selective and structured state space models.
\end{abstract}

\section{From Recurrence to Selection: A Survey of State Space Models in Sequence Modeling}

The fundamental challenge of sequence modeling lies in the persistent tension between capturing long-range dependencies and maintaining computational scalability. Early paradigms, rooted in the Foundations of Sequence Modeling and Recurrent Architectures, provided the initial framework for temporal processing but were often hampered by sequential bottlenecks and gradient instability. While The Transformer Era and Efficient Attention Mechanisms revolutionized the field by enabling massive parallelization and high-capacity learning, the inherent quadratic complexity of self-attention created a significant barrier to modeling increasingly long and complex data streams, necessitating a new class of efficient architectures.

To overcome these scaling constraints, recent research has pivoted toward the Emergence of Structured State Space Models, which offer a path toward linear-time complexity by leveraging continuous-time dynamics. This trajectory has culminated in The Mamba Revolution and Selective State Spaces,

\subsection{Foundations of Sequence Modeling and Recurrent Architectures}

The foundational challenge in sequence modeling is the inherent difficulty of training Recurrent Neural Networks (RNNs), specifically regarding the stability of gradients (\texttt{arXiv:1211.5063}). To address the vanishing and exploding gradient problems identified in earlier literature, researchers have developed both optimization strategies and specialized architectures. For instance, gradient norm clipping serves as a direct mechanism to mitigate exploding gradients (\texttt{arXiv:1211.5063}), while the Adam optimizer (\texttt{arXiv:1412.6980}) provides a computationally efficient first-order gradient-based method. Adam utilizes adaptive estimates of lower-order moments for stochastic objective functions, offering invariance to diagonal rescaling and making it particularly suitable for large-scale, non-stationary, or noisy gradient environments (\texttt{arXiv:1412.6980}).

Architectural evolution has focused on refining how information is gated and preserved. Gated Recurrent Neural Networks, such as LSTMs and GRUs, were introduced to outperform traditional tanh-based recurrent units in applications like speech signal and polyphonic music modeling (\texttt{arXiv:1412.3555}). While these gated units improve information flow, Unitary Evolution RNNs (\texttt{arXiv:1511.06464}) take a more rigorous mathematical approach to long-term dependencies. By employing unitary weight matrices where eigenvalues have an absolute value of exactly 1 in the complex domain, these models circumvent the gradient issues that plague standard RNNs (\texttt{arXiv:1511.06464}). However, to maintain efficiency and avoid the high computational cost of repeated eigendecompositions, these architectures rely on the composition of structured matrices as building blocks (\texttt{arXiv:1511.06464}).

A secondary evolutionary path involves transitioning from purely sequential processing to parallelizable structures to overcome the latency bottlenecks of traditional recurrence. Quasi-Recurrent Neural Networks (QRNNs) (\texttt{arXiv:1611.01576}) attempt to balance these needs by alternating parallel convolutional layers with a minimalist recurrent pooling function. This hybrid approach achieves significant efficiency, providing up to 16 times faster training and testing times compared to standard recurrent models (\texttt{arXiv:1611.01576}). In contrast, Gated Convolutional Networks (\texttt{arXiv:1612.08083}) move entirely away from the recurrent paradigm. By using stacked convolutional networks and a simplified gating mechanism, they allow for parallelization over sequential tokens, reducing latency by an order of magnitude compared to recurrent baselines (\texttt{arXiv:1612.08083}). While QRNNs retain a recurrent element through pooling to maintain sequence modeling capabilities, Gated Convolutional Networks prioritize the speed of non-recurrent, parallelized architectures for tasks such as language modeling (\texttt{arXiv:1611.01576}, \texttt{arXiv:1612.08083}).

\subsection{The Transformer Era and Efficient Attention Mechanisms}

The introduction of the Transformer architecture (\texttt{arXiv:1706.03762}) marked a fundamental shift in sequence modeling by dispensing with recurrence and convolution entirely. By relying solely on multi-head self-attention to draw global dependencies, this architecture achieved increased parallelizability and reduced training time compared to previous recurrent models. This foundation enabled significant evolutionary leaps: BERT (\texttt{arXiv:1810.04805}) introduced deep bidirectional representations by jointly conditioning on both left and right context, a method subsequently refined by RoBERTa (\texttt{arXiv:1907.11692}) through more robust hyperparameter tuning and increased training data. The paradigm was further unified by T5 (\texttt{arXiv:1910.10683}), which converted diverse language tasks into a single text-to-text format, and extended to computer vision via ViT (\texttt{arXiv:2010.11929}), which applies pure Transformers directly to image patches.

Despite these successes, the quadratic complexity of standard self-attention creates a significant bottleneck for long-sequence modeling. To mitigate this, research has diverged into three primary efficiency strategies: sparse attention, kernel-based linear attention, and hardware-aware optimization.

Sparse attention mechanisms attempt to reduce complexity by limiting the attention pattern. Longformer (\texttt{arXiv:2004.05150}) achieves linear scaling by combining local windowed attention with task-motivated global attention, whereas Big Bird (\texttt{arXiv:2007.14062}) utilizes a sparse attention mechanism that maintains the model's status as a universal approximator and Turing complete system.

In contrast to these structural sparsity approaches, other methods focus on approximating the attention mechanism itself to achieve linear complexity. The Linear Transformer (\texttt{arXiv:2006.16236}) uses the associativity property of matrix products via kernel feature maps, which transforms the architecture into an iterative formulation equivalent to an RNN. Similarly, the Performer (\texttt{arXiv:2009.14794}) employs the FAVOR+ mechanism to approximate softmax attention-kernels using positive orthogonal random features. While both approaches aim for \$O(N)\$ complexity, they differ in their mathematical approach to approximating the attention kernel.

A third direction, represented by FlashAttention (\texttt{arXiv:2205.14135}), shifts the focus from mathematical approximation to hardware efficiency. Rather than changing the attention mechanism, FlashAttention uses IO-aware tiling to minimize memory reads and writes between GPU HBM and SRAM, providing fast and memory-efficient exact attention. The proliferation of these distinct methodologies—ranging from sparse patterns and kernel approximations to hardware-level optimizations—necessitated the creation of the Long Range Arena (\texttt{arXiv:2011.04006}), a unified benchmark suite designed to evaluate efficient Transformers across various long-context scenarios.

\subsection{Emergence of Structured State Space Models}

The development of structured state space models (SSMs) has progressed from fundamental memory projection theories to highly efficient, hardware-aware architectures capable of competing with Transformers. This evolution began with the HiPPO framework, which established a mathematical foundation for optimal online compression of continuous signals using polynomial projections [\texttt{arXiv:2008.07669}]. By deriving the Legendre Memory Unit (LMU) from first principles, HiPPO provided a method for recurrent neural networks to approximate functions via projection onto polynomial bases [\texttt{arXiv:2008.07669}].

Building directly upon these foundations, the Structured State Space (S4) model introduced a new parameterization for the state matrix \$A\$ using a Normal Plus Low-Rank (NPLR) structure [\texttt{arXiv:2111.00396}]. This innovation enabled stable diagonalization and efficient computation via a Cauchy kernel, allowing S4 to solve the 16k-length Path-X task with 88\% accuracy—a feat impossible for prior models [\texttt{arXiv:2111.00396}]. While S4 achieved 60× faster generation than Transformers on WikiText-103, its performance was noted to be driven primarily by HiPPO initialization rather than the NPLR parameterization itself, and it remained susceptible to numerical instability if \$A\$ contained eigenvalues in the right half-plane [\texttt{arXiv:2111.00396}].

Subsequent research focused on simplifying these complex structures. The Linear State-Space Layer (LSSL) generalized recurrent, convolutional, and neural differential equation representations by introducing trainable structured matrices \$A\$ [\texttt{arXiv:2110.13985}]. This trend toward simplification was furthered by the discovery that Diagonal State Spaces (DSS) could match the performance of S4's more complex NPLR structure [\texttt{arXiv:2203.14343}]. This led to the development of S4D, which systematically parameterizes diagonal SSMs to recover the S4 kernel, offering improved implementation and computation [\texttt{arXiv:2206.11893}]. Further theoretical refinement by \texttt{arXiv:2206.12037} generalized orthogonal basis projections to explain the S4 model through exponentially-warped Legendre polynomials [\texttt{arXiv:2206.12037}].

Recent advancements have addressed expressivity and architectural efficiency. To bridge the gap between SSMs and Transformers in language modeling, the H3 model introduced a hybrid attention-SSM architecture and the "FlashConv" algorithm for improved hardware utilization [\texttt{arXiv:2212.14052}]. Simultaneously, the transition from single-input, single-output (SISO) to multi-input, multi-output (MIMO) models via parallel scans significantly enhanced efficiency [\texttt{arXiv:2208.04933}]. Finally, while diagonal models like DSS improved efficiency, they were critiqued for their inability to express certain autoregressive processes [\texttt{arXiv:2303.09489}]. This was addressed by the SpaceTime architecture, which utilizes a companion matrix parameterization to enable autoregressive modeling and long-horizon forecasting for time series [\texttt{arXiv:2303.09489}].

\subsection{The Mamba Revolution and Selective State Spaces}

Recent research has focused on bridging the gap between the parallelizable training of Transformers and the efficient inference of Recurrent Neural Networks (RNNs). RWKV (\texttt{arXiv:2305.13048}) addresses this by employing a linear attention mechanism, which allows the model to be formulated as both a Transformer and an RNN, thereby achieving constant computational and memory complexity during inference. A similar objective is pursued by the Linear Recurrent Unit (LRU) (\texttt{arXiv:2303.06349}), which utilizes linearized and diagonalized recurrence to enable parallelizable training. While RWKV relies on linear attention to achieve efficiency, LRU focuses on using signal propagation arguments to recover the performance of deep SSMs in long-range reasoning tasks (\texttt{arXiv:2303.06349}).

The paradigm shifted with the introduction of Mamba (\texttt{arXiv:2312.00752}), which provides a significant improvement over previous subquadratic-time architectures, such as gated convolution and linear attention. The core innovation of Mamba is the Selective State Space Model (SSM), which introduces input-dependent parameters. This selection mechanism allows the model to dynamically decide which information to propagate or forget based on the current token, supported by a hardware-aware parallel algorithm. However, the introduction of input-dependency necessitates careful management of training stability. S7 (\texttt{arXiv:2410.03464}) addresses this by implementing stable reparameterization for input-dependent state transitions. While Mamba provides the mechanism for selection, S7 offers a critical refinement by using reparameterization to control gradient norms, thereby preventing the exploding or vanishing gradient problems that can plague input-dependent models. This makes S7 particularly effective for complex domains like neuromorphic event-based datasets and biological time series (\texttt{arXiv:2410.03464}).

Further refinements to the Mamba architecture are found in the Structured State Space Duality (SSD) framework (\texttt{arXiv:2405.21060}). SSD establishes a formal connection between SSMs and attention variants through the decomposition of structured semiseparable matrices. This refinement of Mamba’s selective SSM is notably more efficient, performing 2-8X faster than the original Mamba (\texttt{arXiv:2405.21060}). In contrast to these dynamic, input-dependent models, Spectral State Space Models (\texttt{arXiv:2312.06837}) offer an alternative approach by using fixed convolutional filters via a spectral filtering algorithm. Unlike the selective mechanisms in Mamba or S7, Spectral SSMs do not require learning input-dependent transitions, making them a distinct choice for predicting long-range dependencies in synthetic dynamical systems (\texttt{arXiv:2312.06837}).

\subsection{Theoretical Analysis and Mathematical Formalisms}

The theoretical landscape of state-space models (SSMs) is characterized by a fundamental tension between approximation capacity and the management of temporal dependencies. A primary limitation of standard SSMs is their capacity constraint, which arises from the absence of nonlinear activation along the temporal direction (\texttt{arXiv:2309.13414}). To resolve this, \texttt{arXiv:2309.13414} demonstrates that stacking SSMs with layer-wise nonlinearity allows them to function as universal approximators for continuous sequence-to-sequence relationships. However, a notable weakness of this approach is that the models still exhibit exponential decaying memory (\texttt{arXiv:2309.13414}).

Different mathematical strategies have been proposed to mitigate this "curse of memory." One approach focuses on the stability of the recurrent structure; \texttt{arXiv:2311.14495} introduces stable reparameterization of recurrent weights to prevent the model from converging to the stability boundary, thereby alleviating the exponential decay limitation. In contrast, rather than focusing on weight stability, \texttt{arXiv:2410.03158} adopts an information-theoretic formalism to address memory efficiency. By utilizing mutual information and rate-distortion theory, this work introduces a selective gating mechanism that dynamically filters the hidden state based on input relevance. This represents a shift from the structural stability focus of \texttt{arXiv:2311.14495} toward a compression-oriented approach that optimizes the trade-off between information retention and memory usage.

The broader mathematical frameworks used to describe these models also diverge. While \texttt{arXiv:2403.16899} provides a control-theoretic overview that connects linear SSMs to modern foundation models and dynamical systems, \texttt{arXiv:2410.03158} relies on information theory to justify selective gating for memory compression. These diverse perspectives are unified by the need for rigorous evaluation; for instance, \texttt{arXiv:2506.05678} proposes a synthetic benchmarking framework using controllable memory functions. This allows for the characterization of temporal dependencies and the comparison of SSMs against other architectures, such as Transformers and RNNs, by generating a continuum of tasks with varying temporal complexity.

\subsection{Advanced Architectures and Specialized Domains}

Recent research has moved beyond standard State Space Model (SSM) implementations to develop specialized architectures that address specific structural or computational limitations. One primary direction is the hybridization of SSMs with other paradigms to bridge performance gaps. For instance, the Block-State Transformer (BST) combines an SSM sublayer for long-range contextualization with a Block Transformer sublayer for short-term representation, achieving a >10x speed increase at the layer level compared to Block-Recurrent Transformers when parallelized (\texttt{arXiv:2306.09539}).

A significant portion of recent innovation focuses on extending SSMs to non-linear or structured data types, such as graphs and spatiotemporal sequences. While earlier approaches might have applied SSM modules to pre-extracted sequences, the Message-Passing State-Space Model (MP-SSM) improves upon this by embedding SSM computation directly into the Message-Passing Neural Network (MPNN) framework, enabling permutation-equivariant, long-range propagation (\texttt{arXiv:2505.18728}). Similarly, GraphSSMs extend SSM theory to temporal graphs by integrating discretized continuous-time linear dynamical systems with a Laplacian regularization term to capture structural information (\texttt{arXiv:2406.00943}). In the spatiotemporal domain, Convolutional SSMs (ConvSSM) utilize parallel scans to achieve subquadratic parallelization, resulting in training that is 3X faster than ConvLSTM and generation 400X faster than Transformers (\texttt{arXiv:2310.19694}).

Another evolutionary trend involves the integration of SSMs with biologically inspired or modular frameworks. To address the limitations of RNN-based Spiking Neural Networks (SNNs), researchers have integrated SSMs with SNNs to facilitate long-range sequence modeling (\texttt{arXiv:2401.00955}). This is further refined in SpikingSSMs, which incorporate dendritic neuron structures to achieve a 90\% average network sparsity through light-weight surrogate dynamic networks (\texttt{arXiv:2408.14909}). While these SNN-based models focus on sparsity and biological realism, SlotSSMs focus on modularity by replacing the monolithic state vector with multiple independent "slots" that interact via a self-attention bottleneck, specifically targeting object-centric learning and 3D visual reasoning (\texttt{arXiv:2406.12272}).

These advancements represent a broader shift toward a taxonomy of recurrent solutions designed to counter the limitations of the Transformer era (\texttt{arXiv:2406.09062}). Whether through the modularity of SlotSSMs (\texttt{arXiv:2406.12272}), the structural integration of MP-SSMs (\texttt{arXiv:2505.18728}), or the efficiency of hybrid BSTs (\texttt{arXiv:2306.09539}), the field is increasingly characterized by tailoring the SSM mathematical formalism to the specific inductive biases of the target domain.

\subsection{The 2025 Frontier: Multi-Scale, Video, and Adaptive SSMs}

The current research landscape for State Space Models (SSMs) in 2025 is characterized by a transition from monolithic architectures toward specialized, multi-resolution, and unified frameworks designed to overcome the inherent limitations of single-scale dependency and computational scaling.

A primary thrust in recent literature is the development of multi-scale and hierarchical architectures to capture complex dependencies. MS-SSM (\texttt{arXiv:2512.23824}) introduces Multi-Resolution Analysis (MRA) via trainable dilated convolutions and scale-dependent \$\textbackslash{}bar\{A\}\$ initialization, allowing the model to process sequence dynamics across multiple resolutions. This approach significantly outperforms single-scale models on hierarchical tasks, with MS-SSM achieving 63.04\% on ListOps compared to Mamba's 38.02\% (\texttt{arXiv:2512.23824}). However, while MS-SSM demonstrates empirical strength, it lacks formal theoretical justification for its scale-dependent initialization and provides only qualitative complexity claims rather than rigorous FLOPs/parameter comparisons (\texttt{arXiv:2512.23824}).

Parallel to multi-scale efforts, researchers are seeking to unify disparate architectures. Prefix-Scannable Models (PSMs) (\texttt{arXiv:2506.10918}) propose a generalization of state aggregation that allows arbitrary functions, such as softmax attention, to unify RNNs and linear transformers, achieving \$O(1)\$ amortized compute per token and \$\textbackslash{}log(N)\$ memory. This unification is furthered by the framework in \texttt{arXiv:2510.09389}, which treats output coefficients as outputs of autonomous linear dynamical systems to reveal common mathematical themes across Transformers, SSMs, and gated linear RNNs.

Adaptivity and selectivity also remain central themes. SeRpEnt (\texttt{arXiv:2501.11729}) extends the Mamba architecture by introducing a selective resampling mechanism that compresses sequences based on information content. In contrast, the Geometric SSM (\texttt{arXiv:2505.17932}) offers an alternative to selective SSMs like Mamba; it achieves selectivity through a dynamic residual generator that excites distinct invariant subspaces, thereby maintaining the Linear Time-Invariant (LTI) property. This represents a significant departure from the time-varying matrices used in Mamba-style models.

Furthermore, the application domain of SSMs is expanding into complex, high-dimensional data. For video modeling, \texttt{arXiv:2505.20171} utilizes a block-wise SSM scanning scheme combined with dense local attention to trade spatial consistency for extended temporal memory, moving beyond simple retrofitting toward true causal video world modeling. In the realm of graph learning, Message-Passing SSMs (MP-SSM) (\texttt{arXiv:2505.18728}) integrate SSM principles into the Message-Passing Neural Network framework to enable permutation-equivariant, long-range information propagation for spatiotemporal graphs. Finally, L2T-Hyena (\texttt{arXiv:2511.05926}) introduces an adaptive "Learn-to-Teach" framework, using a Dynamic Loss Network to adapt the student's loss function online, enhancing the vanilla Hyena architecture.

\subsection{Tensions and Open Problems}

A fundamental tension in sequence modeling concerns the trade-off between the global expressive power of attention mechanisms and the computational efficiency of recurrent or convolutional structures. While the standard Transformer architecture, as established in \texttt{arXiv:1706.03762}, has become the industry benchmark, its quadratic complexity remains a significant bottleneck for long-range dependencies. Efforts to mitigate this, such as the segment-level recurrence introduced in \texttt{arXiv:1901.02860}, attempt to extend context beyond fixed lengths, yet recent surveys like \texttt{arXiv:2412.11211} suggest that even these advanced attention-based models—including sparse variants like the Longformer in \texttt{arXiv:2004.05150} or efficient variants like the Reformer in \texttt{arXiv:2001.04451}—struggle to match the linear-time scaling and memory efficiency offered by State-Space Models (SSMs). This debate is further complicated by the realization that many "efficient" Transformers are essentially approximating recurrent dynamics, a connection highlighted in \texttt{arXiv:2006.16236}, which suggests that the distinction between attention and recurrence may be more fluid than previously assumed.

A second major debate centers on the architectural complexity required to achieve optimal memory and state representation. For several years, the field has grappled with whether structured state-space matrices are necessary for performance. While the HiPPO framework in \texttt{arXiv:2008.07669} emphasizes the importance of optimal polynomial projections for long-term memory, subsequent work in \texttt{arXiv:2203.14343} argues that diagonal state spaces are just as effective as more complex structured versions, potentially simplifying hardware implementation. However, this drive toward simplicity faces pushback: recent critiques, such as \texttt{arXiv:2302.06646}, argue that continuous kernel convolutions may lack the hardware efficiency of simpler alternatives, while \texttt{arXiv:2306.09539} suggests that even hardware-efficient convolutions can be outperformed by block-state transformer architectures. This tension is most visible in the transition from S4 to Mamba; while the selective mechanism in \texttt{arXiv:2312.00752} provides significant gains, recent investigations in \texttt{arXiv:2503.18970} and \texttt{arXiv:2412.11211} question whether these selective SSMs truly resolve the fundamental trade-offs between stability, expressivity, and the ability to handle complex, non-stationary dynamics.

Finally, there is an emerging disagreement regarding the modeling of discrete versus continuous state transitions. Some researchers advocate for the efficiency of discrete state-space models to capture temporal patterns, as seen in \texttt{arXiv:2303.09489}, which suggests that simplicity in discrete states can outperform more complex continuous models. Conversely, others argue that the continuous nature of SSMs is essential for capturing the underlying physics of a sequence, a view that underpins the development of more sophisticated, time-varying dynamics. This conflict highlights an open problem: how to design a model that maintains the mathematical rigor of continuous-time systems while benefiting from the computational advantages of discrete-time hardware acceleration.

\subsection{Outlook}

The frontier of sequence modeling is rapidly shifting from merely extending context length to mastering the complexity of the data itself. We are moving beyond static, linear-time architectures toward models capable of handling multi-scale dynamics (\texttt{arXiv:2512.23824}), geometric constraints, and complex graph-structured dependencies (\texttt{arXiv:2505.18728}). As the field matures, the focus is pivoting toward "world models" that can simulate long-context video environments (\texttt{arXiv:2505.20171}) and adaptive frameworks that allow models to "learn to teach" themselves through coefficient dynamics (\texttt{arXiv:2511.05926}). The ultimate goal is a unified architecture that seamlessly bridges the sequential-parallel duality (\texttt{arXiv:2506.10918}), providing the efficiency of an RNN with the massive-scale training stability of a Transformer.

\section*{References}
\begin{enumerate}\small
  \item On the difficulty of training Recurrent Neural Networks (2012). \texttt{arXiv:1211.5063}
  \item Adam: A Method for Stochastic Optimization (2014). \texttt{arXiv:1412.6980}
  \item Learning Phrase Representations using RNN Encoder-Decoder for
  Statistical Machine Translation (2014). \texttt{arXiv:1406.1078}
  \item Neural Machine Translation by Jointly Learning to Align and Translate (2014). \texttt{arXiv:1409.0473}
  \item Empirical Evaluation of Gated Recurrent Neural Networks on Sequence
  Modeling (2014). \texttt{arXiv:1412.3555}
  \item Deep Residual Learning for Image Recognition (2015). \texttt{arXiv:1512.03385}
  \item Unitary Evolution Recurrent Neural Networks (2015). \texttt{arXiv:1511.06464}
  \item Layer Normalization (2016). \texttt{arXiv:1607.06450}
  \item Gaussian Error Linear Units (GELUs) (2016). \texttt{arXiv:1606.08415}
  \item Pointer Sentinel Mixture Models (2016). \texttt{arXiv:1609.07843}
  \item Language Modeling with Gated Convolutional Networks (2016). \texttt{arXiv:1612.08083}
  \item The LAMBADA dataset: Word prediction requiring a broad discourse context (2016). \texttt{arXiv:1606.06031}
  \item Quasi-Recurrent Neural Networks (2016). \texttt{arXiv:1611.01576}
  \item Attention Is All You Need (2017). \texttt{arXiv:1706.03762}
  \item Decoupled Weight Decay Regularization (2017). \texttt{arXiv:1711.05101}
  \item BERT: Pre-training of Deep Bidirectional Transformers for Language
  Understanding (2018). \texttt{arXiv:1810.04805}
  \item GLUE: A Multi-Task Benchmark and Analysis Platform for Natural Language
  Understanding (2018). \texttt{arXiv:1804.07461}
  \item Think you have Solved Question Answering? Try ARC, the AI2 Reasoning
  Challenge (2018). \texttt{arXiv:1803.05457}
  \item Generating Long Sequences with Sparse Transformers (2019). \texttt{arXiv:1904.10509}
  \item Transformer-XL: Attentive Language Models Beyond a Fixed-Length Context (2019). \texttt{arXiv:1901.02860}
  \item RoBERTa: A Robustly Optimized BERT Pretraining Approach (2019). \texttt{arXiv:1907.11692}
  \item Exploring the Limits of Transfer Learning with a Unified Text-to-Text
  Transformer (2019). \texttt{arXiv:1910.10683}
  \item Compressive Transformers for Long-Range Sequence Modelling (2019). \texttt{arXiv:1911.05507}
  \item Megatron-LM: Training Multi-Billion Parameter Language Models Using
  Model Parallelism (2019). \texttt{arXiv:1909.08053}
  \item HiPPO: Recurrent Memory with Optimal Polynomial Projections (2020). \texttt{arXiv:2008.07669}
  \item Transformers are RNNs: Fast Autoregressive Transformers with Linear
  Attention (2020). \texttt{arXiv:2006.16236}
  \item Language Models are Few-Shot Learners (2020). \texttt{arXiv:2005.14165}
  \item Long Range Arena: A Benchmark for Efficient Transformers (2020). \texttt{arXiv:2011.04006}
  \item Reformer: The Efficient Transformer (2020). \texttt{arXiv:2001.04451}
  \item Longformer: The Long-Document Transformer (2020). \texttt{arXiv:2004.05150}
  \item Rethinking Attention with Performers (2020). \texttt{arXiv:2009.14794}
  \item Linformer: Self-Attention with Linear Complexity (2020). \texttt{arXiv:2006.04768}
  \item An Image is Worth 16x16 Words: Transformers for Image Recognition at
  Scale (2020). \texttt{arXiv:2010.11929}
  \item Big Bird: Transformers for Longer Sequences (2020). \texttt{arXiv:2007.14062}
  \item Efficient Content-Based Sparse Attention with Routing Transformers (2020). \texttt{arXiv:2003.05997}
  \item Efficiently Modeling Long Sequences with Structured State Spaces (2021). \texttt{arXiv:2111.00396}
  \item Combining Recurrent, Convolutional, and Continuous-time Models with
  Linear State-Space Layers (2021). \texttt{arXiv:2110.13985}
  \item CKConv: Continuous Kernel Convolution For Sequential Data (2021). \texttt{arXiv:2102.02611}
  \item On the Parameterization and Initialization of Diagonal State Space
  Models (2022). \texttt{arXiv:2206.11893}
  \item Diagonal State Spaces are as Effective as Structured State Spaces (2022). \texttt{arXiv:2203.14343}
  \item Simplified State Space Layers for Sequence Modeling (2022). \texttt{arXiv:2208.04933}
  \item Hungry Hungry Hippos: Towards Language Modeling with State Space Models (2022). \texttt{arXiv:2212.14052}
  \item FlashAttention: Fast and Memory-Efficient Exact Attention with
  IO-Awareness (2022). \texttt{arXiv:2205.14135}
  \item It's Raw! Audio Generation with State-Space Models (2022). \texttt{arXiv:2202.09729}
  \item How to Train Your HiPPO: State Space Models with Generalized Orthogonal
  Basis Projections (2022). \texttt{arXiv:2206.12037}
  \item Spectral State Space Models (2023). \texttt{arXiv:2312.06837}
  \item Convolutional State Space Models for Long-Range Spatiotemporal Modeling (2023). \texttt{arXiv:2310.19694}
  \item Effectively Modeling Time Series with Simple Discrete State Spaces (2023). \texttt{arXiv:2303.09489}
  \item Block-State Transformers (2023). \texttt{arXiv:2306.09539}
  \item Structured State Space Models for In-Context Reinforcement Learning (2023). \texttt{arXiv:2303.03982}
  \item StableSSM: Alleviating the Curse of Memory in State-space Models through
  Stable Reparameterization (2023). \texttt{arXiv:2311.14495}
  \item Mamba: Linear-Time Sequence Modeling with Selective State Spaces (2023). \texttt{arXiv:2312.00752}
  \item Simple Hardware-Efficient Long Convolutions for Sequence Modeling (2023). \texttt{arXiv:2302.06646}
  \item State-space Models with Layer-wise Nonlinearity are Universal
  Approximators with Exponential Decaying Memory (2023). \texttt{arXiv:2309.13414}
  \item State space models can express n-gram languages (2023). \texttt{arXiv:2306.17184}
  \item Resurrecting Recurrent Neural Networks for Long Sequences (2023). \texttt{arXiv:2303.06349}
  \item RWKV: Reinventing RNNs for the Transformer Era (2023). \texttt{arXiv:2305.13048}
  \item Mathematical Formalism for Memory Compression in Selective State Space
  Models (2024). \texttt{arXiv:2410.03158}
  \item Deep Learning-based Approaches for State Space Models: A Selective
  Review (2024). \texttt{arXiv:2412.11211}
  \item State-Space Modeling in Long Sequence Processing: A Survey on Recurrence in the Transformer Era (2024). \texttt{arXiv:2406.09062}
  \item Learning Long Sequences in Spiking Neural Networks (2024). \texttt{arXiv:2401.00955}
  \item State Space Models as Foundation Models: A Control Theoretic Overview (2024). \texttt{arXiv:2403.16899}
  \item S7: Selective and Simplified State Space Layers for Sequence Modeling (2024). \texttt{arXiv:2410.03464}
  \item State Space Models on Temporal Graphs: A First-Principles Study (2024). \texttt{arXiv:2406.00943}
  \item Slot State Space Models (2024). \texttt{arXiv:2406.12272}
  \item Oscillatory State-Space Models (2024). \texttt{arXiv:2410.03943}
  \item SpikingSSMs: Learning Long Sequences with Sparse and Parallel Spiking
  State Space Models (2024). \texttt{arXiv:2408.14909}
  \item Transformers are SSMs: Generalized Models and Efficient Algorithms
  Through Structured State Space Duality (2024). \texttt{arXiv:2405.21060}
  \item MS-SSM: A Multi-Scale State Space Model for Efficient Sequence Modeling (2025). \texttt{arXiv:2512.23824}
  \item Advancing Intelligent Sequence Modeling: Evolution, Trade-offs, and Applications of State- Space Architectures from S4 to Mamba (2025). \texttt{arXiv:2503.18970}
  \item Message-Passing State-Space Models: Improving Graph Learning with Modern Sequence Modeling (2025). \texttt{arXiv:2505.18728}
  \item Long-Context State-Space Video World Models (2025). \texttt{arXiv:2505.20171}
  \item Geometric SSM: LTI State Space Models for Selective Tasks (2025). \texttt{arXiv:2505.17932}
  \item SeRpEnt: Selective Resampling for Expressive State Space Models (2025). \texttt{arXiv:2501.11729}
  \item Numerical Investigation of Sequence Modeling Theory using Controllable Memory Functions (2025). \texttt{arXiv:2506.05678}
  \item L2T-Hyena: Enhancing State-Space Models with an Adaptive Learn-to-Teach Framework (2025). \texttt{arXiv:2511.05926}
  \item Design Principles for Sequence Models via Coefficient Dynamics (2025). \texttt{arXiv:2510.09389}
  \item Rethinking the long-range dependency in Mamba/SSM and transformer models (2025). \texttt{arXiv:2509.04226}
  \item Sequential-Parallel Duality in Prefix Scannable Models (2025). \texttt{arXiv:2506.10918}
  \item Time-Varying Deep State Space Models for Sequences with Switching Dynamics (2026). \texttt{arXiv:2605.15311}
\end{enumerate}

\end{document}
